\documentclass[10pt,twocolumn]{article}

\usepackage[margin=1in]{geometry}
\usepackage{amsmath,amssymb}
\usepackage{graphicx}
\usepackage{booktabs}
\usepackage{hyperref}
\usepackage{xcolor}
\usepackage{microtype}

\hypersetup{
  colorlinks=true,
  linkcolor=blue!70!black,
  citecolor=blue!70!black,
  urlcolor=blue!70!black,
}

\title{\textbf{Self-Supervised Rectified Flow with Similarity-Based
       Noisy Pairs for Low-Dose CT Denoising}}
\author{Tim Sereda}

\begin{document}
\maketitle

\section{Proposed Method}
\label{sec:method}

We propose a denoiser that uses \emph{no} clean targets, built on rectified
flow, and designed for the spatially correlated noise of FBP-reconstructed
LDCT. The method has two ingredients: an \emph{unconditional} flow whose
learned velocity field is a Noise2Noise estimator, and a
\emph{similarity-based} construction of the noisy training pairs that keeps
the paired noise independent despite spatial correlation.

\subsection{Rectified flow and its Noise2Noise estimator}
Let $s\sim p_{\text{data}}$ be the (unobserved) clean slice and let
$x_0=s+n_0,\ x_1=s+n_1$ be two noisy observations with
$\mathbb{E}[n_i\mid s]=0$ and $n_0\perp n_1\mid s$. Rectified
flow~\cite{lipman2022flow,liu2022flow} defines the linear interpolant
\begin{equation}
  z_t = (1-t)\,x_0 + t\,x_1,\qquad t\in[0,1],
\end{equation}
whose constant velocity is $\dot z_t = x_1-x_0$. We train an
\emph{unconditional} velocity network $v_\theta(z,t)$ (inputs: the
interpolant $z$ and time $t$ only) with the flow-matching loss
\begin{equation}
  \mathcal{L}(\theta)=
  \mathbb{E}_{s,n_0,n_1,\,t\sim\mathcal{U}[0,1]}
  \big\|\,v_\theta(z_t,t)-(x_1-x_0)\,\big\|^2 .
  \label{eq:loss}
\end{equation}

\paragraph{Proposition (denoising via the marginal velocity).}
The minimiser of \eqref{eq:loss} is the conditional expectation
$v^\star(z,t)=\mathbb{E}[\,x_1-x_0\mid z_t=z\,]$. At $t=0$ we have
$z_0=x_0$, so
\begin{equation}
  x_0 + v^\star(x_0,0)
  = \mathbb{E}[x_1\mid x_0]
  = \mathbb{E}[s\mid x_0],
  \label{eq:mmse}
\end{equation}
the MMSE denoiser, because $n_1$ is zero-mean and independent of $x_0$
given $s$. Thus a single velocity evaluation recovers the
Noise2Noise~\cite{lehtinen2018noise2noise} optimum from noisy pairs alone.

\paragraph{Why the flow must be unconditional.}
If the network is additionally given $x_0$ (as in standard conditional flow
matching for denoising), then $(z_t,x_0,t)$ determines the target exactly,
$x_1=(z_t-(1-t)x_0)/t$: the regression target is fully observed and the
network reproduces the \emph{noisy} $x_1$ rather than averaging over it,
destroying \eqref{eq:mmse}. Dropping the conditioning makes each $z_t$
consistent with many pairs, forcing the network to predict their average
velocity, where the independent noise cancels.

\subsection{Similarity-based pair construction}
Equation~\eqref{eq:mmse} requires two observations whose noise is
independent given the signal. FBP reconstruction smears each detector
reading across neighbouring pixels, so spatial neighbours share noise; the
downsampling pairs of Neighbor2Neighbor~\cite{huang2021neighbor2neighbor}
and ZS-N2N~\cite{mansour2023zsn2n} are therefore only valid when the noise
correlation length is below the pixel pitch. We adopt the non-local
construction of Noise2Sim~\cite{niu2020noise2sim}: for each query patch we
retrieve its $K$ most similar patches elsewhere in the slice (smallest
sum-of-squared-differences), \emph{excluding} a radius $r$ around the query
so that matched patches lie beyond the noise correlation length.
Self-similarity makes their underlying signal nearly identical while the
spatial separation makes their noise approximately independent, yielding a
valid pair $(x_0,x_1)$ from a single noisy image. We treat the exclusion
radius $r$ as the key correlated-noise hyper-parameter. As a baseline
variant (\textbf{v1}) we also use spatial-downsampling pairs; the
similarity-based construction is \textbf{v2}.

\subsection{Training and inference}
\textbf{Training.} Per noisy slice: (i)~build endpoints $(x_0,x_1)$ by the
similarity construction; (ii)~sample $t\sim\mathcal{U}[0,1]$ and form $z_t$;
(iii)~take a gradient step on \eqref{eq:loss} with the unconditional
$v_\theta$. No clean targets, no noise model, one shared network across the
dataset.

\noindent\textbf{Inference.} Given a test slice $x$, the one-step estimate
is $\hat s = x + v_\theta(x,0)$. Optionally we refine with a few Euler steps
of the learned ODE $\dot z = v_\theta(z,t)$ initialised at $z_0=x$. We do
\emph{not} integrate to the noisy endpoint distribution: the denoiser is the
posterior mean of \eqref{eq:mmse}, not a transported sample.

\subsection{Hypothesis and ablations}
By \eqref{eq:mmse}, the one-step estimate coincides with similarity-based
Noise2Noise regression. The flow formulation is therefore justified only if
\emph{multi-step refinement improves the structured residual} that one-step
MMSE leaves behind. Our central ablation isolates this: one-step vs.\
$k$-step inference ($k\in\{1,2,4,8\}$); sensitivity to the exclusion radius
$r$; and comparison against direct Noise2Sim regression and the supervised
upper bound. A positive multi-step effect is the contribution; a null effect
is itself a clean finding that narrows the claim to the
correlated-noise-aware pairing.

% -----------------------------------------------------------------------
\bibliographystyle{plain}
\begin{thebibliography}{9}

\bibitem{lipman2022flow}
Y.~Lipman, R.~T.~Q. Chen, H.~Ben-Hamu, M.~Nickel, and M.~Le.
\newblock Flow matching for generative modeling.
\newblock In \emph{International Conference on Learning Representations}, 2023.

\bibitem{liu2022flow}
X.~Liu, C.~Gong, and Q.~Liu.
\newblock Flow straight and fast: Learning to generate and transfer data with
  rectified flow.
\newblock In \emph{International Conference on Learning Representations}, 2023.

\bibitem{lehtinen2018noise2noise}
J.~Lehtinen, J.~Munkberg, J.~Hasselgren, S.~Laine, T.~Karras, M.~Aittala,
  and T.~Aila.
\newblock {Noise2Noise}: Learning image restoration without clean data.
\newblock In \emph{International Conference on Machine Learning}, 2018.

\bibitem{huang2021neighbor2neighbor}
T.~Huang, S.~Li, X.~Jia, H.~Lu, and J.~Liu.
\newblock {Neighbor2Neighbor}: Self-supervised denoising from single noisy
  images.
\newblock In \emph{IEEE Conference on Computer Vision and Pattern Recognition},
  2021.

\bibitem{mansour2023zsn2n}
Y.~Mansour and R.~Heckel.
\newblock Zero-shot {Noise2Noise}: Efficient image denoising without any data.
\newblock In \emph{IEEE Conference on Computer Vision and Pattern Recognition},
  2023.

\bibitem{niu2020noise2sim}
C.~Niu and G.~Wang.
\newblock {Noise2Sim}, similarity-based self-learning for image denoising.
\newblock \emph{arXiv preprint arXiv:2011.03384}, 2020.

\end{thebibliography}

\end{document}
